% Generated by `taskgen course examples` from arithmetic.fft_mul (seed 1).
% Do not edit: change examples.yaml and regenerate.
\begin{example}[умножение с помощью БПФ в $\FF_p$]\label{ex:fftmul-ex}
Вычислить произведение целых чисел \[856 \cdot 990\] с помощью БПФ в поле вычетов $\mathbb{Z}/p\mathbb{Z}$. Обосновать выбор параметров алгоритма.

\par\noindent\textit{Решение.} Цифры (младшие первыми): $a = (6, 5, 8)$, $b = (0, 9, 9)$. Свёртка имеет $3 + 3 - 1 = 5$ коэффициентов, поэтому длина преобразования $N = 8$ (степень двойки). Коэффициенты свёртки не превосходят $81 \cdot \min(3, 3) = 243$, значит нужно $p > 243$ и $N \mid p - 1$ (чтобы существовал корень степени $N$ из единицы): $p = 257$. $g = 3$ --- первообразный корень, $\omega = g^{(p-1)/N} = 3^{32} \equiv 64$, $\omega^{-1} \equiv 253$, $N^{-1} \equiv 225\ (\textrm{mod}\ 257)$. Степени $\omega^j$: $1, 64, 241, 4, 256, 193, 16, 253$.\par $A_k = \sum_j a_j \omega^{jk}$, $B_k = \sum_j b_j \omega^{jk}$, $C_k = A_k B_k$, $c_j = N^{-1} \sum_k C_k \omega^{-jk}$ (всё по модулю $p$):\[\begin{array}{c|rrrrrrrr}k & 0 & 1 & 2 & 3 & 4 & 5 & 6 & 7 \\ \hline A_k & 19 & 198 & 175 & 154 & 9 & 72 & 78 & 114 \\ B_k & 18 & 175 & 104 & 180 & 0 & 51 & 135 & 108 \\ C_k & 85 & 212 & 210 & 221 & 0 & 74 & 250 & 233 \\ \hline c_k & 0 & 54 & 99 & 117 & 72 & 0 & 0 & 0\end{array}\]Переносы: $\sum_j c_j 10^j = 0 \cdot 10^{0} + 54 \cdot 10^{1} + 99 \cdot 10^{2} + 117 \cdot 10^{3} + 72 \cdot 10^{4} + 0 \cdot 10^{5} + 0 \cdot 10^{6} + 0 \cdot 10^{7} = 847440$.\[856 \cdot 990 = 847440.\]
\end{example}
